\documentclass{article}
\usepackage{notes}

\setcounter{section}{0}
\renewcommand{\thesection}{1.\arabic{section}}

\begin{document}
\stdtitle{ANALYSIS II}{Exercise Sheet 1}{Jakob Schildhauer}

\section{Recap Questions}
\begin{enumerate}
    \item No, consider the counterexample $x_n = \frac{(-1)^n + 3}{2}$.
    \item $f(x) = \ln x + C$. 
    \item \dots
\end{enumerate}

\section{Title}
\begin{enumerate}
    \item \dots 
    \item \dots 
    \item \dots
\end{enumerate}

\section{Title}
\begin{enumerate}
    \item \dots 
    \item \dots 
    \item \dots
\end{enumerate}

\section{Title}
\begin{enumerate}
    \item \dots 
    \item \dots 
    \item \dots
\end{enumerate}

\section{Title}
\begin{enumerate}
    \item \dots 
    \item \dots 
    \item \dots
\end{enumerate}

\section{Title}
\begin{enumerate}
    \item \dots 
    \item \dots 
    \item \dots
\end{enumerate}
\end{document}